% Formalización pública del resultado autoral de polarización causal.

\section{Causal polarization field and radial freedom loss}
\label{sec:md-polarizacion-causal-definicion}

Let \((\mathcal M,g)\) be a spherically symmetric geometry written in
Painlevé--Gullstrand coordinates,
\begin{equation}
 \mathrm ds^2
 =-c^2\,\mathrm dt^2
  +\bigl(\mathrm dr+v(r)\,\mathrm dt\bigr)^2
  +r^2\,\mathrm d\Omega_2^2,
 \label{eq:md-polarizacion-metrica-pg}
\end{equation}
where \(v(r)\geq0\) measures the radial inclination of the causal cones with respect
to the selected temporal congruence. The dimensionless field
\begin{equation}
 \mathcal P(r):=\frac{v(r)}{c}
 \label{eq:md-polarizacion-parametro}
\end{equation}
is the intensity of causal polarization. Over the spatial distribution
orthogonal to the temporal congruence, the tensor is defined
\begin{equation}
 \Pi(r)=\mathcal P(r)^2\,\widehat r\otimes\widehat r,
 \label{eq:md-polarizacion-tensor}
\end{equation}
whose principal direction is radial and whose tangential directions
remain degenerate. This definition introduces no absolute direction:
the principal direction is covariantly determined by the gradient of the
areal radius.

The null radial generators of \eqref{eq:md-polarizacion-metrica-pg}
satisfy
\begin{equation}
 \frac{\mathrm dr}{\mathrm dt}=-v(r)\pm c
 =c\bigl(-\mathcal P(r)\pm1\bigr).
 \label{eq:md-polarizacion-generatrices}
\end{equation}
Therefore, the set of radially accessible signs changes cardinality
upon crossing \(\mathcal P=1\).

\begin{theorem}[Mandatory causal direction of fall]
\label{thm:md-polarizacion-caida-obligada}
In the geometry \eqref{eq:md-polarizacion-metrica-pg} the following three exhaustive alternatives are verified:
\begin{enumerate}
 \item if \(0\leq\mathcal P<1\), one future-directed radial generator increases \(r\)
 and the other decreases it;
 \item if \(\mathcal P=1\), the outgoing generator satisfies
 \(\mathrm dr/\mathrm dt=0\) and defines the horizon;
 \item if \(\mathcal P>1\), both future radial generators satisfy
 \(\mathrm dr/\mathrm dt<0\).
\end{enumerate}
In the third regime, the radial degree does not disappear from the manifold, but ceases
to constitute a bidirectional spatial choice: every future-directed causal vector
has a component directed toward smaller areal radii. The interior is, in
this exact sense, a region of compelled causal infall.
\end{theorem}

\begin{proof}
The three conclusions are obtained by comparing the two values
\(-\mathcal P\pm1\) of \eqref{eq:md-polarizacion-generatrices} with zero. For
\(\mathcal P>1\), even the branch with positive sign is negative. The
conclusion concerns the causal cone rather than a particular choice of
temporal coordinate.
\end{proof}

\section{Schwarzschild realization and collective gradient of fall}
\label{sec:md-polarizacion-schwarzschild}

For a total mass \(M\), the Schwarzschild realization in the chart
\eqref{eq:md-polarizacion-metrica-pg} is given by
\begin{equation}
 v(r)=c\sqrt{\frac{R}{r}},
 \qquad
 R=\frac{2GM}{c^2},
 \qquad
 \mathcal P(r)^2=\frac{R}{r}.
 \label{eq:md-polarizacion-schwarzschild}
\end{equation}
The horizon is simultaneously characterized by
\begin{equation}
 r=R,
 \qquad
 \mathcal P=1,
 \qquad
 g^{\mu\nu}\partial_\mu r\,\partial_\nu r=0.
 \label{eq:md-polarizacion-horizonte}
\end{equation}
The radial acceleration of the falling congruence is recovered as the gradient
of the kinetic intensity of the field:
\begin{equation}
 a_r=-\frac12\frac{\mathrm d}{\mathrm dr}v(r)^2
 =-\frac{GM}{r^2}.
 \label{eq:md-polarizacion-gradiente}
\end{equation}
Thus, radial gravitation is expressed as the collective gradient of infall and
not as a force added after the trajectories have been fixed.

\begin{corollary}[Causal Macropolarization of the Gravitational Interior]
\label{cor:md-macropolarizacion-interior}
In the realization \eqref{eq:md-polarizacion-schwarzschild}, the region
\(r<R\) satisfies \(\mathcal P>1\). The black hole constitutes to
causal macropolarization: the principal direction of \(\Pi\) exceeds the
unit threshold and aligns both future radial families toward the decrease of the
areal radius.
\end{corollary}

\section{Correspondencia with the \texorpdfstring{\(3\)}{3} regimes of the TRIT}
\label{sec:md-polarizacion-trit}

The geometric reader of the TRIT is implemented through
\begin{equation}
 \tau(r)=\operatorname{sgn}\!\bigl(1-\mathcal P(r)^2\bigr)
 \in\{+1,0,-1\}.
 \label{eq:md-polarizacion-lector-trit}
\end{equation}
The correspondence is
\begin{equation}
 \begin{array}{c|c|c|c}
 \tau & \mathcal P & \text{geometric regime} &
 \text{future radial orientation}\\ \hline
 +1 & \mathcal P<1 & \text{exterior} & \text{input and output}\\
 0  & \mathcal P=1 & \text{horizon} & \text{outgoing null generator}\\
 -1 & \mathcal P>1 & \text{interior} & \text{necessary descent}
 \end{array}
 \label{eq:md-polarizacion-tabla-trit}
\end{equation}
The temporal involution of TRIT exchanges the exterior and interior regimes
and fixes the horizon. The change \(+1\leftrightarrow-1\) does not eliminate the memory
of the route: it modifies its causal orientation and preserves the boundary
\(\tau=0\) as to gluing surface.

\Needspace{30\baselineskip}
\section{Local closure of matter and collective gravitational alignment}
\label{sec:md-polarizacion-materia}

Let \(\gamma_a\) be a persistent material section with radial projector
\(\widehat r_a\otimes\widehat r_a\) and closure weight \(w_a\geq0\). For a
finite set of sections, define the order parameter
\begin{equation}
 \Pi_N
 =\frac{1}{\sum_{a=1}^{N}w_a}
  \sum_{a=1}^{N}w_a\,
  \widehat r_a\otimes\widehat r_a.
 \label{eq:md-polarizacion-colectiva}
\end{equation}
An individual particle realizes a persistent local closure of a route. The
gravitational interior appears when the collective field fixes a principal
causal direction for all future routes. Consequently, two
mathematically distinct objects are not identified:
\begin{equation}
\begin{aligned}
 \text{particle}
 &=\text{persistent material section},\\
 \text{black hole}
 &=\text{macroscopic causal alignment above the threshold}.
\end{aligned}
 \label{eq:md-polarizacion-distincion}
\end{equation}

The complete causal sequence of this residence is fixed by
\begin{equation}
\begin{aligned}
 \text{APP--TRIT--TPK state}
 &\longrightarrow
 \text{material section}
 \longrightarrow
 \text{polarization field }\Pi\\
 &\longrightarrow
 \mathcal P=1
 \longrightarrow
 \text{horizon}
 \longrightarrow
 \mathcal P>1
 \longrightarrow
 \text{necessary causal descent}.
\end{aligned}
 \label{eq:md-polarizacion-cadena}
\end{equation}
The following chapter prolongs this structure toward sheet inversion,
the advanced--retarded purification, Kruskal geometry, and horizon
information.

\section{Geometric conditions of causal macropolarization: areal radius, cone tilt, horizon, and interior falloff}
\label{sec:md-polarizacion-control}

The result simultaneously requires that the radius used be areal, that \(v(r)\) tilt the cones according to \eqref{eq:md-polarizacion-metrica-pg}, that the horizon satisfy \(\mathcal P=1\), and that the two slopes in \eqref{eq:md-polarizacion-generatrices} be negative in the interior. These identities distinguish the effective loss of radial freedom from a merely dimensional metaphor and provide direct control of causal macropolarization.
